\documentclass[10pt,a4paper]{amsart}
\usepackage{amsmath,amssymb,amsthm,graphicx}
\usepackage{verbatim}
\newenvironment{reference}{\comment}{\endcomment}
%\newenvironment{reference}{}{}
\newenvironment{slogan}{\comment}{\endcomment}
\newenvironment{history}{\comment}{\endcomment}
% For commutative diagrams we use Xy-pic
\usepackage[all]{xy}

% We use 2cell for 2-commutative diagrams.
\xyoption{2cell}
% Theorem environments.
%
\theoremstyle{plain}
\newtheorem{theorem}[subsection]{Theorem}
\newtheorem{proposition}[subsection]{Proposition}
\newtheorem{lemma}[subsection]{Lemma}

\theoremstyle{definition}
\newtheorem{definition}[subsection]{Definition}
\newtheorem{example}[subsection]{Example}
\newtheorem{exercise}[subsection]{Exercise}
\newtheorem{situation}[subsection]{Situation}

\theoremstyle{remark}
\newtheorem{remark}[subsection]{Remark}
\newtheorem{remarks}[subsection]{Remarks}

\numberwithin{equation}{subsection}

% Macros
%
\def\lim{\mathop{\mathrm{lim}}\nolimits}
\def\colim{\mathop{\mathrm{colim}}\nolimits}
\def\Spec{\mathop{\mathrm{Spec}}}
\def\Hom{\mathop{\mathrm{Hom}}\nolimits}
\def\Ext{\mathop{\mathrm{Ext}}\nolimits}
\def\SheafHom{\mathop{\mathcal{H}\!\mathit{om}}\nolimits}
\def\SheafExt{\mathop{\mathcal{E}\!\mathit{xt}}\nolimits}
\def\Sch{\mathit{Sch}}
\def\Mor{\mathop{\mathrm{Mor}}\nolimits}
\def\Ob{\mathop{\mathrm{Ob}}\nolimits}
\def\Sh{\mathop{\mathit{Sh}}\nolimits}
\def\NL{\mathop{N\!L}\nolimits}
\def\CH{\mathop{\mathrm{CH}}\nolimits}
\def\proetale{{pro\text{-}\acute{e}tale}}
\def\etale{{\acute{e}tale}}
\def\QCoh{\mathit{QCoh}}
\def\Ker{\mathop{\mathrm{Ker}}}
\def\Im{\mathop{\mathrm{Im}}}
\def\Coker{\mathop{\mathrm{Coker}}}
\def\Coim{\mathop{\mathrm{Coim}}}

% Boxtimes
%
\DeclareMathSymbol{\boxtimes}{\mathbin}{AMSa}{"02}

%
% Macros for moduli stacks/spaces
%
\def\QCohstack{\mathcal{QC}\!\mathit{oh}}
\def\Cohstack{\mathcal{C}\!\mathit{oh}}
\def\Spacesstack{\mathcal{S}\!\mathit{paces}}
\def\Quotfunctor{\mathrm{Quot}}
\def\Hilbfunctor{\mathrm{Hilb}}
\def\Curvesstack{\mathcal{C}\!\mathit{urves}}
\def\Polarizedstack{\mathcal{P}\!\mathit{olarized}}
\def\Complexesstack{\mathcal{C}\!\mathit{omplexes}}
% \Pic is the operator that assigns to X its picard group, usage \Pic(X)
% \Picardstack_{X/B} denotes the Picard stack of X over B
% \Picardfunctor_{X/B} denotes the Picard functor of X over B
\def\Pic{\mathop{\mathrm{Pic}}\nolimits}
\def\Picardstack{\mathcal{P}\!\mathit{ic}}
\def\Picardfunctor{\mathrm{Pic}}
\def\Deformationcategory{\mathcal{D}\!\mathit{ef}}
\title{Direct and inverse image along a closed immersion reflects finite generation of modules}
\author{The Stacks Project Authors}
\date{stacks-project:modules.tex:local-snapshot-20261012; GNU FDL 1.2 or later}
\begin{document}
\maketitle
\noindent\textit{Editorial note.} All delivered mathematical text is the author's own, copied from the pinned Stacks Project chapter source sources/stacks-modules/pinned-original-source.tex; no statement or proof step was rewritten, reordered or removed.\par\medskip
\noindent\textit{Editorial note.} Preamble provenance: the Stacks Project keeps one single shared preamble for every chapter; this card takes it directly from the staged sources/stacks-modules/preamble.tex as three verbatim spans (lines 8-12, 14-18 and 154-217), so no derived file was written and nothing was added or redefined. The wrapper class and the base packages are supplied by the assembler.\par\medskip
\noindent\textit{Editorial note.} Bibliography provenance: neither the counted statement nor its proof contains a citation (this lane rescans the whole span before writing the card), so the card carries no bibliography and no entry was reconstructed.\par\medskip
\noindent\textit{Editorial note.} Selection provenance: this card was opened under the file-based control contract (CONTROL.md), with the mechanical pre-card gate run on its statement and proof interval before the card was written.\par\medskip
\noindent\textit{Allowed standard prerequisites (editorial).}\begin{itemize}
\item The notion of an \textbackslash{}mathcal\{O\}\_X-module of finite type: every point has a neighbourhood on which finitely many sections generate the module. The author's section on modules of finite type, in which the counted statement occurs. Formulates both directions of the equivalence.
\item The description of the stalk of a direct image sheaf and the surjectivity of \textbackslash{}mathcal\{O\}\_X \textbackslash{}to i\_*\textbackslash{}mathcal\{O\}\_Z at the points of the closed subset. Used inside the counted proof at source lines 2140-2142. Lets the proof conclude surjectivity of the displayed map from its behaviour on stalks.
\end{itemize}\medskip
\begin{lemma}
\label{lemma-i-star-reflects-finite-type}
Let $i : (Z, \mathcal{O}_Z) \to (X, \mathcal{O}_X)$ be a morphism of ringed
spaces. Assume $i$ is a homeomorphism onto a closed subset of $X$ and
that $\mathcal{O}_X \to i_*\mathcal{O}_Z$ is surjective.
Let $\mathcal{F}$ be an $\mathcal{O}_Z$-module.
Then $i_*\mathcal{F}$ is of finite type if and only if
$\mathcal{F}$ is of finite type.
\end{lemma}

\begin{proof}
Suppose that $\mathcal{F}$ is of finite type.
Pick $x \in X$. If $x \not \in Z$, then $i_*\mathcal{F}$
is zero in a neighbourhood of $x$ and hence finitely generated in
a neighbourhood of $x$. If $x = i(z)$, then choose an open neighbourhood
$z \in V \subset Z$ and sections $s_1, \ldots, s_n \in \mathcal{F}(V)$
which generate $\mathcal{F}$ over $V$. Write $V = Z \cap U$ for some open
$U \subset X$. Note that $U$ is a neighbourhood of $x$. Clearly the
sections $s_i$ give sections $s_i$ of $i_*\mathcal{F}$ over $U$.
The resulting map
$$
\bigoplus\nolimits_{i = 1, \ldots, n} \mathcal{O}_U
\longrightarrow
i_*\mathcal{F}|_U
$$
is surjective by inspection of what it does on stalks
(here we use that $\mathcal{O}_X \to i_*\mathcal{O}_Z$
is surjective). Hence $i_*\mathcal{F}$
is of finite type.

\medskip\noindent
Conversely, suppose that $i_*\mathcal{F}$ is of finite type.
Choose $z \in Z$. Set $x = i(z)$. By assumption there exists
an open neighbourhood $U \subset X$ of $x$, and sections
$s_1, \ldots, s_n \in (i_*\mathcal{F})(U)$ which generate $i_*\mathcal{F}$
over $U$. Set $V = Z \cap U$. By definition of $i_*$ the sections
$s_i$ correspond to sections $s_i$ of $\mathcal{F}$ over $V$.
The resulting map
$$
\bigoplus\nolimits_{i = 1, \ldots, n} \mathcal{O}_V
\longrightarrow
\mathcal{F}|_V
$$
is surjective by inspection of what it does on stalks. Hence
$\mathcal{F}$ is of finite type.
\end{proof}

\end{document}
